% =====================================================================================
% sections/S2_methods.tex -- Supplementary Section S2: further details of the methods.
% All `% src:` paths are relative to the root of the repository.
% =====================================================================================
\section{Methods: further details}\label{S2:sec}

\subsection{The fixed point sets of the heat problem}

\begin{figure}[tbp]
\centering
\includegraphics[width=\textwidth]{s3_methods_points.pdf}
\caption{The four fixed interior point sets for problem \PH\ on the space--time
square, $\Nr=256$: cell-centred and node-centred $16\times16$ grids, and the
\strat{random} and \strat{Sobol} sets of seed 0 as generated by the benchmark code. Grey
dots: the 300 initial and boundary points, the same in every panel. Shaded: the strip
$0<t<1/32$ between the initial line and the first column of the cell-centred grid. The
strategies \strat{resample} and \strat{RAD} change their sets during training
(Table~\ref{s3_methods:tab:samplers}) and are not shown.}
\label{s3_methods:fig:points}
\end{figure}
% src: code/s3_methods_points.py -> figures/s3_methods_points.pdf, data/s3_methods_points.json

Figure~\ref{s3_methods:fig:points} shows the four fixed sets for \PH. The cell-centred grid
leaves the strip $0<t<1/32$ next to the initial line without a residual point; the
node-centred grid has a column on the initial line but its next column is at $t=1/15$.

\subsection{Implementation of the low-dimensional runs}\label{S2:sec:impl}

% src: research_benchmark/results/runs/*.json (timing.lbfgs_batched_fevals / 1500 = 7.28, 2.51, 2.59, 2.15, 1.49, 1.84)
The multi-seed runs use our own L-BFGS for a stack of networks, with Armijo backtracking
(Section~\ref{sC_details:sec:settings}); a stack advances only when its slowest member has
accepted a step, and the stacks used between 1.5 and 7.3 evaluations per iteration, an
upper bound for each member. Section~\ref{sC_details:sec:validation} compares the two L-BFGS
implementations.

% src: research_benchmark/results/runs/*.json (timing.batch_size 25 or 50, device mps); research_benchmark/scripts/run_budget_sweep.py (specs: 3 samplers x 5 seeds = stacks of 15, device mps); data/s4_lowdim_gridcontrol.json (stacks: batch_size 40 and 20, device mps); data/s5_hard_laplace3d_smooth.json (note: run_batched on CPU, 5 seeds stacked; device cpu); notes/VERIFICATION.md (sec. 2.5, caveats 4 and 5); data/checks/machine_load.json; research_highdim/scripts/hd_core.py (torch.set_num_threads(1))
All low-dimensional experiments are run by one PyTorch package (\code{pinnbench}); the
benchmark trains 25 or 50 networks per stack, independent in exact arithmetic, stacked along
a batch axis on the GPU of an Apple-silicon machine (budget sweep of
Section~\ref{S3:sec:budget}: stacks of 15; control experiment of
Section~\ref{s4_lowdim:sec:strategy}: stacks of 40 and 20 on the GPU; \PLthreeS: 5 networks on
the CPU). The driver is validated against single-network training in
Section~\ref{sC_details:sec:validation}; in floating point the result of a network depends on
the size of its stack (Section~\ref{sC_details:sec:repro}). The $d$-dimensional runs use one
CPU thread. The machine, with 12 cores, was shared with other jobs throughout, which is why
Sections~\ref{s5_hard:sec} and~\ref{s6_highdim:sec} quote more than one time measurement
wherever a cost is stated.

\subsection{The penalty bias for \texorpdfstring{\Pone}{LapD} at \texorpdfstring{$d=2$}{d = 2}}

\begin{observation}[Size of the penalty bias for \Pone\ at $d=2$]\label{s3_methods:obs:robin}
% src: verify_research_highdim/results/v_robin_fd.json (rel_l2_N200; N100 vs N200 differ by <= 2.1e-4 relative)
% src: verify_research_highdim/scripts/v_robin_fd.py
For \Pone\ in dimension $d=2$, where $\uex=x_1x_2$ and $f=0$, the Robin problem of
Proposition~\ref{s3_methods:prop:robin}(a) was solved by second-order finite differences
(five-point scheme with ghost points; the solutions on grids of 100 and 200 cells per
side differ by less than 0.03 per cent in the quantity reported). The relative $L^2$
distance $\norm{\ubeta-\uex}/\norm{\uex}$ is $0.288$, $\sci{5.00}{-2}$,
$\sci{5.48}{-3}$ and $\sci{5.54}{-4}$ for $\bet=1$, $10$, $100$ and $1000$, that is,
between $0.50/\bet$ and $0.55/\bet$ for $\bet\ge10$, in line with part~(c); the bound of
part~(c) and the sharper bound of Corollary~\ref{s6_highdim:cor:bound}(b) are compared with these
values in Observation~\ref{sA_proofs:obs:sharp}, and the bias in dimension $d=3$ to $20$ is computed
in Observation~\ref{sA_proofs:obs:biasd} and, exactly, in Proposition~\ref{sA_proofs:prop:exact}.
% src: verify_research_highdim/results/v_summary.txt (rows 'robin'); notes/VERIFICATION.md (sec. 3.1, V5 and V9)
Deep Ritz networks trained with the protocol of Section~\ref{s3_methods:sec:highdim} (one
run for each $\bet$, 4000 iterations, with the verification code \textsf{V-dim} of
Section~\ref{s3_methods:sec:verif}) end at $\errL=0.290$ for $\bet=1$ and
$\sci{4.78}{-2}$ for $\bet=10$, that is, within 5 per cent of the bias of the Robin
minimiser and far from zero; at $\sci{1.02}{-2}$ for $\bet=100$, above the bias; and at
$\sci{5.89}{-2}$ for $\bet=1000$, about a hundred times the bias, where the error is that
of the optimisation. The distance of the networks from $\ubeta$ itself was not measured.
\end{observation}

% src: verify_research_highdim/results/v_summary.txt (robin, beta 10); verify_research_highdim/results/v_robin_fd.json (10: rel_l2_N200)
\noindent The bias is the error of the limit of the training, not a lower bound on the
error of a trained network: a network at distance $\norm{v-\ubeta}$ from the minimiser has
an error between $\norm{\ubeta-\uex}-\norm{v-\ubeta}$ and
$\norm{\ubeta-\uex}+\norm{v-\ubeta}$, and the run with $\bet=10$ indeed ends 4.5 per cent
below the bias. What a small $\bet$ does is to make training converge to a function that is
not the solution; in these runs a large $\bet$ in turn slowed the optimisation.
Section~\ref{s6_highdim:sec} returns to both effects. The proof of
Proposition~\ref{s3_methods:prop:robin}(a) also gives the identity
$\Eritz(v)-\Eritz(\ubeta)=\tfrac12\norm{\grad(v-\ubeta)}^2+\bet\,\nbd{v-\ubeta}^2$ quoted in
Section~\ref{s3_methods:sec:ritz}; the minimum of the loss is negative for \Ptwo\ and positive
for \Pone, where $f=0$.

\subsection{Two ways of computing the Laplacian}\label{S2:sec:lap}

The Laplacian in the PINN residual can be computed in two ways. Nested reverse-mode
differentiation takes one further backward pass for each of the $d$ coordinates.
Alternatively the triple $(u,\grad u,\lap u)$ can be propagated forward
through the layers~\citep{li2024forward}. Both are exact: at $d=2$ and $d=10$ the two
versions return the same single-precision loss value, digit for digit as printed (against
the separately written nested version of \textsf{V-dim} the single-precision losses agree
to seven or more significant digits up to $d=20$), and in double precision the Laplacians
differ by less than $10^{-16}$; they differ only in cost (Section~\ref{s6_highdim:sec:cost}).
The main runs of Section~\ref{s6_highdim:sec} use the forward version; five runs of the
weight sweep, made earlier, and \textsf{V-dim} use the nested one.
% src: research_highdim/results/check_core.txt (forward-vs-nested); verify_research_highdim/results/v_check.txt (V1); notes/VERIFICATION.md (sec. 3.1, V1 and V10; sec. 3.4); data/checks/study_records.json (highdim.sweep_runs_before_changes: five sweep runs with nested autograd)
